\documentclass[10pt,a4paper]{amsart}
\usepackage[margin=19mm]{geometry}
\usepackage{amsmath,amssymb,mathtools}
\usepackage[scr=rsfs,cal=euler]{mathalfa}
\usepackage{xfrac}
\usepackage[unicode,hidelinks,bookmarks=false]{hyperref}
\newcommand{\ie}{i.e.\ }
\newcommand{\cf}{cf.\ }
\newcommand{\resp}{respectively\ }
\newcommand{\Subsection}{\subsection}
\providecommand{\nobreakdash}{\nobreak\hbox{-}}
\setlength{\emergencystretch}{2em}
\multlinegap0pt
\usepackage{bm}
\usepackage{bbm}
\usepackage{dsfont}
\usepackage{xspace}
\usepackage{cancel}
\usepackage{mathtools}
\usepackage{amsthm}
\makeatletter
\@ifundefined{thm}{\newtheorem{thm}{Thm}}{}
\@ifundefined{prop}{\newtheorem{prop}[thm]{Proposition}}{}
\makeatother
\title[Picard-Lefschetz theory and alien calculus: a case study]{Picard-Lefschetz theory and alien calculus: a case study}
\author{Si Li, Yong Li, Xinxing Tang}
\date{Source: arXiv:2605.08867v2 (CC BY 4.0)}
\begin{document}
\maketitle

\noindent\emph{Source and licence.} Picard-Lefschetz theory and alien calculus: a case study. Version arXiv:2605.08867v2 (CC BY 4.0), distributed under the Creative Commons Attribution 4.0 International licence, as recorded in the arXiv licence metadata for this exact version and shown on the abstract page https://arxiv.org/abs/2605.08867v2. Full rights notice: licenses/arxiv-2605.08867-CC-BY-4.0.txt.\par\smallskip

\noindent\textit{Editorial scope.} Counted unit: the Prop whose source label is \texttt{prop:gamma-generic-endpoints-W} in the article (source0.tex lines 3639--3688, source label \texttt{prop:gamma-generic-endpoints-W}), delivered together with the article's own complete original proof of it (source0.tex lines 3690--4012, one \texttt{proof} environment of 323 lines). No mathematics is written, repaired or re-derived: statement and proof are the article's own text line for line. Required context retained uncounted: none. Every definition, display and label of those lines is retained unchanged. Omitted from this package and not redistributed: 1--3638, 3689--3689, 4013--5241. Nothing inside a retained excerpt is omitted or altered: every retained body is delivered exactly as the article's lines are written, so no omission receipt of any kind accompanies this package. Citations: the retained excerpts carry no citation command at all, so no bibliography entry of the article is transcribed and no reference is redistributed. Reference adaptation: none was needed and none was made; every reference inside the delivered unit resolves inside it, so no unresolved reference or citation is produced. Printed slips kept literal: no printed word, symbol, hypothesis or number of the counted statement or of its proof was altered. Layout and class: the article's own class is \texttt{amsart} with packages including graphicx, framed, ulem, pgf, tikz, pgfplots, float, mathrsfs, mathtools, mathpple, xy, amsmath, tikz-cd, mathdots; the wrapper is the lane's pinned \texttt{amsart} editorial wrapper at 10pt on A4, which restates the theorem-like environments the retained text needs and adds the article's own macro definitions that the retained text needs (none), with extra wrapper packages (bm, bbm, dsfont, xspace, cancel, mathtools). The delivered numbering is the unit's own. Assets: no retained span contains a figure environment, a graphics-inclusion command, a PDF-page inclusion, a table or a TikZ picture; no image file is copied into this package, the package declares no asset dependency and no image file is redistributed with it. Licence: version arXiv:2605.08867v2 of this article is distributed under the Creative Commons Attribution 4.0 International licence, as recorded in the arXiv licence metadata for this exact version and shown on the abstract page https://arxiv.org/abs/2605.08867v2. A full rights notice is delivered at licenses/arxiv-2605.08867-CC-BY-4.0.txt. Characters: every retained excerpt is pure ASCII, so the delivered engine, XeLaTeX with its default Latin Modern text font, reports no missing character.
\begin{prop}
\label{prop:gamma-generic-endpoints-W}
Let
\[
\theta_-:=\frac{\pi}{2}-\delta,
\qquad
\theta_+:=\frac{\pi}{2}+\delta,
\qquad
0<\delta\ll1.
\]
For \(p_n=2\pi i n\), the stable and unstable manifolds have the following ends.
\begin{enumerate}
\item For \(\theta=\theta_-\), the asymptotics are given by
\[
W^s_{\theta_-}(p_n):
\quad
\begin{cases}
x\to-\infty,\quad y=x\tan\theta_-+\tan\theta_-+2\pi n+o(1),\\[0.1cm]
x\to+\infty,\quad y=\theta_-+2\pi n+O(xe^{-x}),
\end{cases}
\]
and
\[
W^u_{\theta_-}(p_n):
\quad
\begin{cases}
x\to+\infty,\quad y=\theta_- -\pi+2\pi n+O(xe^{-x}),\\[0.1cm]
x\to+\infty,\quad y=\theta_- +\pi+2\pi n+O(xe^{-x}).
\end{cases}
\]
\item For \(\theta=\theta_+\), the asymptotics are given by
\[
W^s_{\theta_+}(p_n):
\quad
\begin{cases}
x\to+\infty,\quad y=\theta_+ +2\pi n+O(xe^{-x}),\\[0.1cm]
x\to+\infty,\quad y=\theta_+ +2\pi(n-1)+O(xe^{-x}),
\end{cases}
\]
and
\[
W^u_{\theta_+}(p_n):
\quad
\begin{cases}
x\to-\infty,\quad y=x\tan\theta_+ +\tan\theta_+ +2\pi n+o(1),\\[0.1cm]
x\to+\infty,\quad y=\theta_+-\pi+2\pi n+O(xe^{-x}).
\end{cases}
\]
\end{enumerate}
\end{prop}

\begin{proof} 
\noindent\textbf{The case $x\rightarrow+\infty$}:

Fix \(n\in\mathbb Z\).  Along the stable and unstable manifolds of
\(p_n=2\pi i n\), the function \(G_\theta\) is constant and equal to
\(G_\theta(p_n)\).  Hence the relevant level equation is
\begin{equation}\label{eq:gamma-level-right-proof}
H_{\theta,n}(x,y):=
G_\theta(x,y)-G_\theta(p_n)
=
e^x\sin(y-\theta)+(x+1)\sin\theta-(y-2\pi n)\cos\theta
=0 .
\end{equation}



By its initial direction, the zero levels set enters the region
\[
  R_{n,\theta}
  :=
  \{\,x>0,\ 2\pi n<y<2\pi n+\theta\,\}.
\]
\begin{itemize}
    \item On the lower horizontal boundary \(y=2\pi n\), for our $\theta$, we have
\[
  H_{\theta,n}(x,2\pi n)
  =
  \sin\theta\,(1+x-e^x)<0,\quad x>0.
\]
    \item On the upper horizontal boundary \(y=2\pi n+\theta\), we have
\[
  H_{\theta,n}(x,2\pi n+\theta)
  =
  (x+1)\sin\theta-\theta\cos\theta .
\]
If \(\theta=\theta_+\), then \(\cos\theta<0\), so the right-hand side is
strictly positive. If \(\theta=\theta_-\), then \(\cos\theta>0\), but
\[
  \sin\theta-\theta\cos\theta>0,\quad 0<\theta<\pi/2.
\] 
Hence
\[
  H_{\theta,n}(x,2\pi n+\theta)>0,
  \qquad x\geq 0 .
\]
\item On the  vertical side \(x=0\), it is easy to check 
\[
  H_{\theta,n}(0,y)>0,
  \qquad 2\pi n<y<2\pi n+\theta.
\]
\end{itemize}
By the above analysis, the right-up branch
cannot leave \(R_{n,\theta}\), otherwise it would have to cross to one of the above boundary pieces. In particular \(y=O(1)\)
along that right end. Hence
\eqref{eq:gamma-level-right-proof} gives
\[
e^x\sin(y-\theta)
=
-(x+1)\sin\theta+(y-2\pi n)\cos\theta
=
O(x).
\]
Thus
\[
\sin(y-\theta)=O(xe^{-x}).
\]
Consequently, on any such right end, \(y-\theta\) must approach a zero of
\(\sin\), and hence there exists \(k\in\mathbb Z\) such that
\[
y=\theta+k\pi+\varepsilon(x),
\qquad
\varepsilon(x)\to 0
\qquad
(x\to+\infty).
\]

Substituting this form into \eqref{eq:gamma-level-right-proof}, we obtain
\[
e^x\sin(k\pi+\varepsilon)
+(x+1)\sin\theta
-(\theta+k\pi+\varepsilon-2\pi n)\cos\theta
=0.
\]
Since
\[
\sin(k\pi+\varepsilon)=(-1)^k\sin\varepsilon
=
(-1)^k\left(\varepsilon+O(\varepsilon^3)\right),
\]
the equation becomes
\[
\bigl((-1)^k e^x-\cos\theta\bigr)\varepsilon
+
(x+1)\sin\theta
-
(\theta+k\pi-2\pi n)\cos\theta
+
O(e^x\varepsilon^3)
=0.
\]
The preliminary estimate \(\sin(y-\theta)=O(xe^{-x})\) gives
\[
\varepsilon(x)=O(xe^{-x}),
\]
and hence
\[
e^x\varepsilon^3=O(x^3e^{-2x}).
\]

Solving the previous equation for \(\varepsilon\), and using
\[
\frac{1}{(-1)^k e^x-\cos\theta}
=
(-1)^k e^{-x}+O(e^{-2x}),
\]
we obtain
\begin{equation*}\label{eq:gamma-right-epsilon-proof}
\varepsilon(x)
=
(-1)^k e^{-x}
\left(
-(x+1)\sin\theta
+
(\theta+k\pi-2\pi n)\cos\theta
\right)
+
O(x^2e^{-2x}).
\end{equation*}
Therefore every right end of \(H_{\theta,n}=0\) in the \(k\)-th horizontal
strip has the asymptotic expansion
\begin{equation}\label{eq:gamma-right-end-proof}
y
=
\theta+k\pi
+
(-1)^k e^{-x}
\left(
-(x+1)\sin\theta
+
(\theta+k\pi-2\pi n)\cos\theta
\right)
+
O(x^2e^{-2x}).
\end{equation}
In particular,
\[
y=\theta+k\pi+O(xe^{-x}),
\qquad x\to+\infty.
\]

It remains to identify whether such a right end is stable or unstable.  Recall that we have
\[
F_\theta(x,y)
=
e^x\cos(y-\theta)-x\cos\theta-y\sin\theta.
\]
Since along \eqref{eq:gamma-right-end-proof}, \(y-\theta=k\pi+\varepsilon(x)\) and \(\varepsilon(x)=O(xe^{-x})\),
\[
\cos(y-\theta)
=
\cos(k\pi+\varepsilon)
=
(-1)^k+O(\varepsilon^2).
\]
Therefore
\[
F_\theta(x,y)
=
(-1)^k e^x+O(x).
\]
Hence, if \(k\) is even, then
\[
F_\theta(x,y)\to+\infty,
\]
so this is a stable, rapid-decay end.  If \(k\) is odd, then
\[
F_\theta(x,y)\to-\infty,
\]
so this is an unstable end.

Thus the stable right ends are precisely the even horizontal asymptotic lines
\[
y=\theta+2\pi \ell+O(xe^{-x}),
\qquad \ell\in\mathbb Z,
\]
whereas the unstable right ends are precisely the odd horizontal asymptotic
lines
\[
y=\theta+(2\ell+1)\pi+O(xe^{-x}),
\qquad \ell\in\mathbb Z.
\]
Specializing these alternatives to the two adjacent strips containing the
branches issuing from \(p_n\) gives exactly the \(x\to+\infty\) endpoint
statements in this proposition.

\medskip
\noindent\textbf{The case $x\rightarrow-\infty$:}


Fix the critical point \(p_n=2\pi i n\).  As before, the stable and unstable
manifolds through \(p_n\) are contained in the level set \eqref{eq:gamma-level-right-proof}.

We now consider the left end, namely an end on which \(x\to-\infty\).  Since
\(\theta=\theta_\pm=\frac{\pi}{2}\pm\delta\) with \(0<\delta\ll1\), we have
\(\cos\theta\neq0\).  Solving \eqref{eq:gamma-level-right-proof} for \(y\), we get
\begin{equation*}\label{eq:gamma-left-solved-proof}
y
=
x\tan\theta
+
\tan\theta
+
2\pi n
+
\frac{e^x}{\cos\theta}\sin(y-\theta).
\end{equation*}
The last term is \(O(e^x)\), because \(\sin(y-\theta)\) is bounded and
\(\theta\) is fixed away from \(\frac{\pi}{2}\) by the small nonzero parameter
\(\delta\).  Therefore
\begin{equation}\label{eq:gamma-left-asymptotic-proof}
y
=
x\tan\theta
+
\tan\theta
+
2\pi n
+
O(e^x),
\qquad x\to-\infty.
\end{equation}
In particular,
\[
y
=
x\tan\theta
+
\tan\theta
+
2\pi n
+
o(1),
\qquad x\to-\infty.
\]
This proves the claimed left-end asymptotic line.

For completeness, let us also record the corresponding asymptotic behavior of
\(F_\theta\), since this determines whether the left end is stable or unstable.
Along the left end, \eqref{eq:gamma-left-asymptotic-proof} gives
\[
y
=
x\tan\theta
+
\tan\theta
+
2\pi n
+
O(e^x).
\]
Substituting this into
\[
F_\theta(x,y)
=
e^x\cos(y-\theta)
-
x\cos\theta
-
y\sin\theta,
\]
we obtain
\[
\begin{aligned}
F_\theta(x,y)
&=
e^x\cos(y-\theta)
-
x\cos\theta
-
\left(
x\tan\theta+\tan\theta+2\pi n+O(e^x)
\right)\sin\theta  \\
&=
-
x\left(
\cos\theta+\frac{\sin^2\theta}{\cos\theta}
\right)
+O(1)=
-\frac{x}{\cos\theta}+O(1).
\end{aligned}
\]
Hence, as \(x\to-\infty\),
\[
F_\theta(x,y)\to+\infty
\quad\text{if}\quad
\cos\theta>0,
\]
whereas
\[
F_\theta(x,y)\to-\infty
\quad\text{if}\quad
\cos\theta<0.
\]
Since \(\cos\theta_->0\) and \(\cos\theta_+<0\), the left end is a stable end
for \(\theta=\theta_-\), and an unstable end for \(\theta=\theta_+\).

Therefore the unique left asymptotic line of the level set through \(p_n\) is
\[
x\to-\infty,
\qquad
y
=
x\tan\theta
+
\tan\theta
+
2\pi n
+
o(1),
\]
and its stable/unstable type is determined by the sign of \(\cos\theta\), as
claimed in this proposition.
\end{proof}

\end{document}
